\section{Data for Statistics and non-analog Methods} \label{sec2.9}

\medskip

\textbf{General remarks}

\medskip

The statistical performance (FOM, ``figure of merit",
see equation \ref{eq1.13} in section \ref{sec1.2.1}) of an EIRENE
run (as for any Monte Carlo application in general) is very
sensitive to the non-analog methods used by EIRENE. Therefore the
input parameters in this block, which define the non-analog setting of an EIRENE run, should only be used, if the user
has carefully worked through the code and has a detailed knowledge
of the Monte Carlo techniques activated by setting the flags in
this block. 

Only the ``Cards for Standard Deviation", which allow to obtain graphical and printed output
regarding statistical standard deviation for all tallies estimated by EIRENE, can be used
with no risk, except that evaluations of variances and covariances
generally tend to slow down the code a bit (few percent) in case
of large meshes.

Hence, as a rule, one
should generally try to find the optimal setting of the non-analog
sampling
flags in this block and estimate the FOM from test runs,
and then use this setting, but without evaluation of standard
deviations, in production runs.

The input flag NLANA in block 1 (section \ref{sec2.1}) de-activates all non-analog sampling
distributions. This option should, for example, be used to derive an
intuitive but physically correct picture from plots of selected
particle trajectories (input block 11, section \ref{sec2.11}),
which otherwise, e.g. in case of suppression of absorption,  may be grossly misleading.
\medskip

\textbf{The Input Block}

\medskip

{\sl*** 9. Data for Statistics and non-analog Methods}

{\small
\begin{verbatim}
 (NLPRCA(IATM)  IATM=1,NATM) , (NLPRCM(IMOL)   IMOL =1,NMOL) ,
 (NLPRCI(IION)  IION=1,NION) , (NLPRCPH(IPHOT) IPHOT=1,NPHOT)
 NPRCSF
 next: read NPRCSF integers (FORMAT 12I6) IPRSF(J) J=1,NPRCSF
 MAXLEV  MAXRAD  MAXPOL  MAXTOR  MAXADD
 DO IN=1,MAXRAD
   READ (IUNIN,66665) ID,NSSPL(IN),PRMSPL(IN)
 ENDDO
 DO IN=1,MAXPOL
   READ (IUNIN,66665) ID,NSSPL(N1ST+IN),PRMSPL(N1ST+IN)
 ENDDO
 DO IN=1,MAXTOR
   READ (IUNIN,66665) ID,NSSPL(N1ST+N2ND+IN),PRMSPL(N1ST+N2ND+IN)
 ENDDO
 DO IN=1,MAXADD
   READ (IUNIN,66665) ID,NSSPL(N1ST+N2ND+N3RD+IN),
                        PRMSPL(N1ST+N2ND+N3RD+IN)
 ENDDO
 WMINV  WMINS  WMINC WMINL
 SPLPAR
\end{verbatim}
}

 {\sl * Cards for Standard Deviation}

\begin{verbatim}
 NSIGVI  NSIGSI  NSIGCI  NSIGI_BGK  NSIGI_COP  NSIGI_SPC
 DO  93  J=1,NSIGVI
    IGH  IIH
 93 CONTINUE
 DO  94  J=1,NSIGSI
    IGHW  IIHW
 94 CONTINUE
 DO  95  J=1,NSIGCI
    IGHC(1,J) IIHC(1,J) IGHC(2,J) IIHC(2,J)
 95 CONTINUE
\end{verbatim}

\medskip

\textbf{Meaning of the Input Variables for Statistics and non-analog
Methods}

\medskip

\begin{list}{ }{}
\item[\textbf{NLPRCA }] conditional expectation estimator (eq.
\ref{eq1.12}) is used for atom species IATM
\item[\textbf{NLPRCM }]
conditional expectation estimator is used for molecule species
IMOL
\item[\textbf{NLPRCI }] conditional expectation estimator is
used for test ion species IION (not ready to use)
\item[\textbf{NLPRCPH }] conditional expectation estimator is
used for photon species IPHOT (in versions 2004 and younger)
\item[\textbf{IPRSF
}] conditional expectation estimator is used, if trajectory points
towards additional surface IPRSF. IPRSF $\leq$ NLIMI, the total number of additional
surfaces read in input block 3B.

NPRCSF surfaces have that property of ``attracting trajectories".
\end{list}
The next lines define so called Splitting and Russian Roulette
surfaces (S$\&$R-surfaces). Any test-particle crossing a
S$\&$R-surface in the negative direction is split into SPLPAR
identical but independent particles with properly reduced weight.
If a particle crosses a S$\&$R-surface in the positive direction,
then it is either killed or continues it's flight with properly
increased weight. The surface orientation can be reversed by using
a negative value for the the splitting parameter SPLPAR.
\begin{list}{ }{}
\item[\textbf{MAXLEV }]
Maximum number of levels for splitting ($\le$ 15)
\item[\textbf{MAXRAD }]
Total number of radial splitting surfaces (-NR1ST $\le$ MAXRAD $\le$ NR1ST)
\begin{list}{ }{}
\item MAXRAD $<$ 0

-MAXRAD is used, the position of the
radial splitting surfaces is automatically defined, and a constant
splitting parameter (SPLPAR, see below) is used for radial splitting and RR.
\item MAXRAD $>$ 0

radial surfaces with numbers NSSPL(IN), IN=1,MAXRAD are
S$\&$R-surfaces. The splitting parameter for surface NSSPL(IN) is
PRMSPL(IN).
\end{list}
\item[\textbf{MAXPOL }] Total number of poloidal splitting surfaces
(0 $\le$ MAXPOL $\le$ NP2ND). The S$\&$R-surfaces and splitting
parameters are selected as in the case of radial surfaces, see
above. 
\item[\textbf{MAXTOR }] Total number of toroidal splitting
surfaces (0 $\le$ MAXTOR $\le$ NT3RD). The S$\&$R-surfaces and
splitting parameters are selected as in the case of radial
surfaces, see above. 
\item[\textbf{MAXADD }] Total number of
additional splitting surfaces (0 $\le$ MAXADD $\le$ NLIMI). The
S$\&$R-surfaces and splitting parameters are selected as in the
case of radial surfaces, see above. 
\item[\textbf{WMINV }] minimum
weight used for suppression of absorption at collisions
(``survival biassing"). If a particle goes into a collision with
weight less than WMINV, then suppression of absorption or any
other non-analog weight correction is abandoned, and the analog
game is played. WMINV acts only for events in the volume including
volume source birth events, but not for events at surfaces.
\item[\textbf{WMINS }] Same as WMINV, but for surface events
(including surface source birth events). 
\item[\textbf{WMINC }]
minimum acceptable weight for conditional expectation estimators
(by abuse of language). More precisely, WMINC is the minimal
acceptable probability for a test flight to reach a particular
cell without collision. If this probability is smaller than WMINC,
the particle track is stopped and restarted. E.g. for WMINC $\ge
1$, the estimator used in the NIMBUS code results (ref.
\cite{kn:Cupini}), whereas for WMINC = 0 each particle path is
integrated according to equation \ref{eq1.12}, until the nearest
non transparent surface along the track is reached, regardless of
any collisions. Periodicity surfaces are regarded as ``transparent" in this context.


Note: strictly speaking this is not a non-analog method, but rather a
particular choice of an unbiased estimator. Hence: the flag
NLANA in input block 1 does not affect flags for conditional
expectation estimators.
\item[\textbf{WMINL }]
to be written

\item[\textbf{SPLPAR }]
splitting parameter for default radial splitting option (MAXRAD $<0$)
\item[\textbf{NSIGVI }]
number of standard deviation profiles for volume averaged tallies
to be estimated.
NSIGVI must be less than or equal to the parameter NSD in PARMUSR
\ref{sec3.1}.
\item[\textbf{NSIGSI }]
number of standard deviation profiles for surface averaged tallies
to be estimated.
NSIGSI must be less than or equal to the parameter NSDW in PARMUSR
\ref{sec3.1}.
\item[\textbf{NSIGCI }]
number of correlation coefficients between volume tallies
to be estimated.
NSIGCI must be less than or equal to the parameter NCV in PARMUSR
\ref{sec3.1}.
\item[\textbf{NSIGI\_BGK }]
if \verb|NSIGI_BGK| $>$ 0, then the standard deviations are
evaluated for all tallies needed for the iteration procedure for
the non-linear BGK collision terms. See subroutine
\verb|STATIS_BGK| in code segment BGK.F
\item[\textbf{NSIGI\_COP }]
if \verb|NSIGI_COP| $>$ 0, then the standard deviations are
evaluated for all tallies needed for the iteration procedure for
the coupling to a plasma fluid model. The relevant tallies are
selected in subroutine \verb|STATIS_COP| in code segment
\verb|COUPLE_....F|, section \ref{sec4.2}
\item[\textbf{NSIGI\_SPC }]
if \verb|NSIGI_SPC| $>$ 0, then the standard deviations are
evaluated for all surface flux spectra defined in sub-block 10F below.
\item[\textbf{IGH }]
Index of species for selected volume averaged tally
\item[\textbf{IIH }]
Index of volume averaged tally for which empirical standard
deviation is to be calculated (see table \ref{tab2}).
\item[\textbf{IGHW }]
Index of species for selected surface averaged tally
\item[\textbf{IIHW }]
Index of surface averaged tally for which empirical standard
deviation is to be calculated (see table \ref{tab3})
\item[\textbf{IGHC1 IIHC1 IGHC2 IIHC2 }]~

Species and tally index for first and second tally, respectively,
between which the correlation coefficient is evaluated
\end{list}

Note: On printout the NSIGVI and NSIGSI standard deviations are
printed as relative standard deviations, in percent. The standard
deviations for the two tallies involved in the NSIGCI arrays are
printed as absolute tallies, i.e., with the same units as the two
tallies themselves.
